\statement{B.83}{Proved}{A whole actual filling kernel in a source subgroup enlarged by the product-one meridian}
Consider the ordered surface and its actual fillings
\[
\begin{split}
Y={}&\{(v,w):v,w\in\mathbb C^*\setminus\mu _5,\
                    v\ne w,\ vw\notin\mu _5\},\\
Y_c={}&\{(v,w):v\in\mathbb C^*\setminus\{1\},\
                 w\in\mathbb C^*\setminus\mu _5,\
                    v\ne w,\ vw\notin\mu _5\},\\
Y_b={}&\{(v,w):v,w\in\mathbb C^*\setminus\{1\},\
                    v\ne w,\ vw\notin\mu _5\}.
\end{split}
\tag{B.83.1}
\]
Their source curves are the corresponding open subsets of \(v=2-w\).
Use the common basepoint \((v_*,w_*)=(399/200,1/200)\).
Let \(J_c\), respectively I, be the actual image of the source
kernel for \(Y\to Y_c\), respectively \(Y\to Y_b\).

There are two specified constant-family source meridians \(c_1,c_4\),
three specified old source lifts \(g_1,g_2,g_5\), and an actual
product-one meridian b, all marked below. Put
\[
 H_*=\langle c_1,c_4,g_1,g_2,g_5\rangle,\qquad
 T=\langle H_*,b\rangle,\qquad
 J_*=\langle\!\langle c_1,c_4\rangle\!\rangle_{H_*}.
\tag{B.83.2}
\]
Then, as actual subgroups of \(\pi _1(Y)\),
\[
\begin{gathered}
 J_*\subseteq J_c,\qquad
 T\cap\ker\bigl(\pi _1(Y)\to\pi _1(Y_c)\bigr)=J_*,\\
 T\cap J_c=J_*,\qquad T/J_*\simeq\mathbb Z\times F_3.
\end{gathered}
\tag{B.83.3}
\]
For the both-family filling put
\[
 I_*=\langle\!\langle c_1,c_4,g_2,g_5\rangle\!\rangle_{H_*}.
\tag{B.83.4}
\]
Then
\[
\begin{gathered}
 I_*\subseteq I,\qquad
 T\cap\ker\bigl(\pi _1(Y)\to\pi _1(Y_b)\bigr)=I_*,\\
 T\cap I=I_*,\qquad T/I_*\simeq\mathbb Z^2.
\end{gathered}
\tag{B.83.5}
\]
The same intersection and quotient statements hold for the actual
closed images in full completions, with \(\mathbb Z,F_3\) replaced
by \(\widehat{\mathbb Z},\widehat F_3\). No induced full topology
on the abstract subgroup T is assumed. This does not identify
\(J_c\) or I with the entire ambient filling kernel.

\begin{proof}
\emph{Actual source marking and the exact action table.}\par
Project to w and use the fiber coordinate \(z=wv\). The twelve
fiber points and the source point are
\[
 0;\quad wu\ (u\in\mu _5);\quad w^2;\quad
 u\ (u\in\mu _5);\qquad z_S=w(2-w).
\tag{B.83.6}
\]
The old base has punctures \(0,\mu _5\). The five values
\(-\mu _5\) are filled transverse fiber nodes. The source has
thirteen additional punctures: 2, the four \(2-u\) with
\(u\ne1\), and the eight \(1\pm\sqrt{1-u}\) with \(u\ne1\).
Its actual group is therefore \(F_{19}\). Filling the constant
family removes the four \(2-u\) punctures; filling both families
also removes the four nontrivial fifth roots from the w-base.
These are the actual source maps \(F_{19}\to F_{15}\) and
\(F_{19}\to F_{11}\).

Write \(\zeta=e^{2\pi i/5}\). The based source meridians at
\(2-\zeta^j\) are denoted \(c_j\), in the actual orientation
of the fiber homomorphism. The source lifts \(g_1,g_2,g_5\)
lie over the distinguished old base lassos at
\(1,\zeta,\zeta^4\), respectively. Their projections are three
members of a free basis of the old \(F_6\) base. The actual
local marking at \(w=1\), transported along the common source
whisker, has retained meridians a,d,b around \(v=1\), \(v=w\),
and \(vw=1\), respectively.

The common paths start at \(1/200\), use the trunk to
\(1/200+i/100\), and straight whiskers to northeast corners
of critical squares of half-side \(10^{-4}\). The exact rational
centers and all words are specified in the certificate cited below.
The initial label order is
\([9,10,8,3,4,2,0,6,5,1,12,11,7]\), with source label12
at position11. The source-adapted fiber basis is
\[
 z_0,c_1,c_2,c_3,c_4,p_1,p_2,p_3,p_4,a,d,b.
\tag{B.83.7}
\]
Its first nine entries are actual source words. Both substitutions
are checked on all twelve generators. Thus these are actual
elements, rather than unbased peripheral classes.

Let \(\gamma_i\) be conjugation by \(g_i\) on the fiber image.
The actual based restrictions are
\[
\begin{split}
 \gamma_2|_{\langle b,c_4\rangle}
        &=\operatorname{Inn}(c_4^{-1}b),\\
 \gamma_5|_{\langle b,c_1\rangle}
        &=\operatorname{Inn}(bc_1^{-1}),\qquad
 \gamma_1(b)=b.
\end{split}
\tag{B.83.8}
\]
The following table gives \(k_\gamma\) in the literal equation
\(\gamma(b)=k_\gamma b k_\gamma^{-1}\):
\[
\begin{array}{c|cccccc}
\gamma&\gamma_1&\gamma_1^{-1}&\gamma_2&\gamma_2^{-1}
             &\gamma_5&\gamma_5^{-1}\\ \hline
k_\gamma&1&1&c_4^{-1}&b^{-1}c_4b
             &bc_1^{-1}b^{-1}&c_1.
\end{array}
\tag{B.83.9}
\]
These identities hold in the actual free generic fiber before
the five node fillings, hence in \(\pi _1(Y)\). They are verified
by the faithful Artin action and the inverse distinguished-meridian
conjugator on all twelve kernel generators. No normalized local
operator is substituted for an actual source operator.

The coefficient comparison is also certified. The rational fifth-root
coordinates differ from the true coordinates by less than \(10^{-30}\).
On all121 segments, \(|w|<6\) and pair separation exceeds
\(10^{-12}\). Linear coefficient deformation changes a pair
difference by less than \(24\cdot10^{-30}\), so it is collision-free
and preserves the initial marking. True extra quadratic critical
values lie in the assigned disks of radius \(10^{-25}\). Thus
the computation represents the actual cyclotomic pencil.

\emph{The initial kernel and the valid stage induction.}\par
The subgroup \(J_*\) is the actual image of the selected source
kernel \(F_5\to F_3\) that kills \(c_1,c_4\). The three surviving
generators project to a free basis subset of the old base; hence
their images remain free in the filled source. In particular
\(J_*\subseteq J_c\), and \(J_*\) is invariant under all three
\(\gamma_i\) and their inverses.

Let \(K_0\) be the image of
\(\ker(F(b,c_1,c_4)\to\langle b\rangle)\). We first show
\(K_0\subseteq J_*\). For \(c=c_1\), the subgroup generated
by its \(\gamma_5\)-orbit is invariant under conjugation by
\(bc^{-1}\), by (B.83.8), and under c, because it contains c.
It is therefore b-normal and contains every \(b^ncb^{-n}\).
For \(c=c_4\), use \(\gamma_2\) and \(b=c(c^{-1}b)\).
These two families generate the entire displayed free-group
kernel. Thus \(K_0\subseteq J_*\), and \(K_0\) is b-normal.

Put \(\Gamma=\langle\gamma_1,\gamma_2,\gamma_5\rangle\), and
let \(K_n\) be generated by the images of \(K_0\) under operator
words of length at most n. Suppose \(K_n\) is b-normal.
For any generator \(\gamma\) or its inverse, (B.83.9) gives
\(\gamma^{-1}(b)=kbk^{-1}\) with \(k\in K_0\subseteq K_n\).
For \(w\in K_n\),
\[
 b\gamma(w)b^{-1}
   =\gamma\bigl(kb(k^{-1}wk)b^{-1}k^{-1}\bigr)
                     \in\gamma(K_n).
\tag{B.83.10}
\]
Using \(b^{-1}\) as well proves that \(\gamma(K_n)\) is
b-normal. Hence \(K_{n+1}\) is b-normal. The essential
hypothesis is \(k\in K_n\); knowing only that k normalizes
the final subgroup would not justify the induction.

The union of these stages is \(J_*\): it is contained in
\(J_*\), contains both source meridians, and is invariant
under the three source generators; conversely \(K_0\) itself
was generated by their source orbits. Thus b normalizes the
whole \(J_*\), including mixed operator orbits.

\emph{The quotient and an actual detector through the filling.}\par
Modulo \(J_*\), table (B.83.9) makes b commute with all three
source lifts. Therefore there is a surjection
\[
                 \mathbb Z\times F_3\longrightarrow T/J_*.
\tag{B.83.11}
\]
The rational function
\[
                    f_b(v,w)=\frac{vw-1}{(w-1)^2}
\tag{B.83.12}
\]
maps both Y and \(Y_c\) to \(\mathbb C^*\). On the actual
source \(v=2-w\), it is identically \(-1\). Define \(\chi_b\) as minus the conventional positive winding
of \(f_b\). The actual inverse-conjugator marking sends a
positive local product-one loop to \(b^{-1}\); hence this
character kills the whole source image and satisfies
\(\chi_b(b)=1\).
Together with the actual w-base projection it identifies the
images of b and the three source lifts with the standard factors
of \(\mathbb Z\times F_3\). The selected base subgroup is a
free factor of \(F_6\); compose the base map with its retraction
when viewing the detector as a map to \(F_3\).
This detector gives a left inverse to
(B.83.11), proving the quotient isomorphism. Both detector
maps factor through \(Y_c\); any element of T killed there
already lies in \(J_*\). The reverse inclusion is the actual
source filling. This proves (B.83.3).

\emph{The both-family consequence.}\par
The loops \(g_2,g_5\) are source meridians of the moving-family
divisors \(w=\zeta,\zeta^4\), filled in \(Y_b\). Hence
\(I_*\subseteq I\). The image of \(H_*\) in
\(T/J_*=\mathbb Z\times F_3\) is exactly \(\{0\}\times F_3\):
its character \(\chi_b\) vanishes and its base image is all
of this free factor. Consequently \(I_*/J_*\) is exactly
the normal closure of \(g_2,g_5\) in that \(F_3\). It is
normal in the whole product, and the quotient is
\(\mathbb Z\times\langle g_1\rangle\).
The character \(\chi_b\) still factors through \(Y_b\).
Winding of \(w-1\), oriented to send \(g_1\) to one, supplies
a second character there. Their product detects this
\(\mathbb Z^2\) quotient and proves (B.83.5).
The equality with \(T\cap I\) uses \(I\) contained in the
ambient filling kernel; no global equality is assumed.

\emph{Actual compact images.}\par
Let \(\overline H_*,\overline T,\overline J_*,\overline I_*\)
be the actual closed images in the full completions.
Here the completed \(J_c\) and I likewise mean the closed images
of the completed full source kernels \( \widehat F_{19}\to
\widehat F_{15}\) and \(\widehat F_{19}\to\widehat F_{11}\).
They are the closures of the discrete images: these actual
source fillings kill subsets of a free basis, so their closed
normal kernels are the closures of their discrete normal kernels.
The completed selected source kernels are closed normal
closures of the listed generators; their images are the
closures of the discrete images above. Continuity preserves
b-normality and all commutation relations. We obtain
continuous surjections
\[
 \widehat{\mathbb Z}\times\widehat F_3
       \longrightarrow\overline T/\overline J_*,
 \qquad
 \widehat{\mathbb Z}^{\,2}
       \longrightarrow\overline T/\overline I_*.
\tag{B.83.13}
\]
The completed actual base maps and winding characters give
left inverses. The selected free-base inclusion is an actual
free-factor inclusion, so its full topology is preserved.
Both surjections are isomorphisms. These same detectors
factor through the completed target fillings, proving the
compact kernel equalities directly. We do not identify
\(\overline T\) with the abstract completion of T.

The actual braid words and source basis are in
\texttt{work\_DE/next\_marginal\_source/}\allowbreak
\texttt{w\_pencil\_words.json} and
\texttt{w\_source\_adapted\_basis.json}, with producers and
genuine outputs. The independent checks
\texttt{work\_CE/code/verify\_w\_pencil\_independent.py} and
\texttt{check\_w\_normalizer\_limitations.py} verify the marking,
all132 kernel-conjugation comparisons, all24 local transports,
the two restrictions, and all inverse conjugators.
The remaining source meridians \(c_2,c_3\) and other old
operators are not covered by the induction. This is a
whole-kernel theorem for T, not for the full ambient group
or the entire source image.
\end{proof}
